\begin{statementbox}
\begin{center}\begin{tikzpicture}[font=\small,>=stealth,line join=round]
\draw[gray!40,rounded corners=3pt] (0,0) rectangle (14.2,8.3);\node at (7.1,7.700000000000001) {周期2的余弦：轴每隔1出现，中心也每隔1出现};
\begin{scope}[shift={(2.3,3.8222222222222224)},x=1.2cm,y=1.277777777777778cm]\draw[->,gray] (-1,0)--(4,0) node[right]{$x$};\draw[->,gray] (0,-1.7)--(0,1.9) node[above]{$y$};\draw[blue,thick,solid,domain=-0.7:3.7,samples=180] plot(\x,{cos(deg(pi*\x))});\fill[orange!85!black] (0.5,0) circle(2.2pt);\fill[orange!85!black] (1.5,0) circle(2.2pt);\fill[orange!85!black] (2.5,0) circle(2.2pt);\fill[orange!85!black] (3.5,0) circle(2.2pt);\foreach \a in {0,1,2,3}{\draw[violet,dashed] (\a,-1.3)--(\a,1.3);\node[above] at (\a,1.35){$\a$};}\end{scope}\node[anchor=west,align=left,text width=5.0cm] at (8.4,4.0) {$f=\cos\pi x,\ T_0=2$\\[9pt]紫线：轴 $x=k$\\[9pt]橙点：中心 $(\frac12+k,0)$\\[9pt]同类间隔 $\frac{T_0}2=1$\\[9pt]但 $f(0)=1\ne f(1)=-1$};
\node[text width=12.6cm,align=center] at (7.1,.65) {半周期是同类对称位置的间隔，不表示每个函数值都在半周期后重复。};
\end{tikzpicture}\end{center}


\textbf{\textcolor{CStatement}{这条结论解决什么问题？}}

有些题只给出“函数隔一段距离又重复”“图像有一条对称轴”或“图像有一个对称中心”，却要求寻找\textbf{其他对称位置、函数周期或函数值}。这时不用先求出函数的表达式，也可以直接推理。

\medskip
\end{statementbox}\clearpage\begin{statementbox}\textbf{\textcolor{CStatement}{先分清三句话}}（以下函数均定义在整个实数范围内）
\begin{itemize}[leftmargin=*,itemsep=0.35em]
\item \textbf{隔 $T$ 又重复：} $T>0$，且对任意实数 $x$ 都有 $f(x+T)=f(x)$。这叫“$T$ 是一个周期”。
\item \textbf{沿 $x=a$ 折过去完全重合：} 对任意 $t$ 都有 $f(a+t)=f(a-t)$。这说明 $x=a$ 是对称轴。
\item \textbf{绕 $(b,c)$ 转半圈完全重合：} 对任意 $t$ 都有 $f(b+t)+f(b-t)=2c$。这说明 $(b,c)$ 是对称中心。
\end{itemize}

\medskip
\textbf{\textcolor{CStatement}{四条常用结论}}
\begin{description}[style=nextline,leftmargin=2.8em,labelsep=0.5em,itemsep=0.55em]
\item[\textbf{1. 一个周期 + 一条对称轴 $\Rightarrow$ 更多对称轴：}]
若 $T$ 是一个周期，$x=a$ 是一条对称轴，那么
\[
\boxed{x=a+\frac{kT}{2}\quad(k=0,\pm1,\pm2,\ldots)}
\]
都是对称轴。\textbf{注意：这些一定有，但不一定只有这些。}

\item[\textbf{2. 一个周期 + 一个对称中心 $\Rightarrow$ 更多对称中心：}]
若 $T$ 是一个周期，$(b,c)$ 是一个对称中心，那么
\[
\boxed{\left(b+\frac{kT}{2},\,c\right)\quad(k=0,\pm1,\pm2,\ldots)}
\]
都是对称中心。\textbf{只移动横坐标，纵坐标仍是 $c$。}

\item[\textbf{3. 两条不同的对称轴 $\Rightarrow$ 一个周期：}]
若 $x=a_1,x=a_2$ 都是对称轴，且 $a_1\ne a_2$，那么
\[
\boxed{P=2|a_2-a_1|}
\]
一定是函数的一个正周期。\textbf{不能直接说它最小。}

\item[\textbf{4. 已知最小正周期，才能列出\emph{全部}对称轴：}]
如果 $T_0$ 是最小正周期，且已知对称轴 $x=a$，那么\textbf{全部}对称轴恰好是
\[
\boxed{x=a+\frac{kT_0}{2}\quad(k\in\mathbb Z).}
\]
\end{description}

\medskip
\end{statementbox}\clearpage\begin{statementbox}
\begin{center}\begin{tikzpicture}[font=\small,>=stealth,line join=round]
\draw[gray!40,rounded corners=3pt] (0,0) rectangle (14.2,8.3);\node at (7.1,7.700000000000001) {轴到中心不一定最近：也可能跨3个四分之一周期};
\begin{scope}[shift={(2.158823529411765,3.8222222222222224)},x=1.7647058823529411cm,y=1.277777777777778cm]\draw[->,gray] (-0.6,0)--(2.8,0) node[right]{$x$};\draw[->,gray] (0,-1.7)--(0,1.9) node[above]{$y$};\draw[blue,thick,solid,domain=-0.3:2.5,samples=180] plot(\x,{cos(deg(pi*\x))});\fill[orange!85!black] (0,1) circle(2.2pt);\fill[orange!85!black] (0.5,0) circle(2.2pt);\fill[orange!85!black] (1,-1) circle(2.2pt);\fill[orange!85!black] (1.5,0) circle(2.2pt);\draw[violet,dashed] (0,-1.3)--(0,1.3);\foreach \v/\L in {.5/\frac12,1/1,1.5/\frac32,2/2}{\node[below=4pt] at (\v,-1.35){$\L$};}\end{scope}\node[anchor=west,align=left,text width=5.0cm] at (8.4,4.0) {$T_0=2$，轴 $x=0$\\[9pt]中心 $(\frac12,0),(\frac32,0)$\\[9pt]$\frac12=T_0/4$\\[9pt]$\frac32=3T_0/4$\\[9pt]一般为奇数个 $T_0/4$};
\node[text width=12.6cm,align=center] at (7.1,.65) {四分之一周期关系使用最小正周期；不能把较远中心也误当最近的中心。};
\end{tikzpicture}\end{center}

\textbf{\textcolor{CStatement}{值得多记的一条：对称轴 + 对称中心}} 

若 $x=a$ 是对称轴，$(b,c)$ 是对称中心，则对任意 $x$，
\[
\boxed{f\bigl(x+2(b-a)\bigr)=2c-f(x).}
\]
意思是：向右（或向左）走 $2|b-a|$，函数值会以 $c$ 为中间高度“上下对换”。再走同样的距离，才回到原值。

如果 $a\ne b$，那么 $4|b-a|$ 是一个正周期；如果 $a=b$，函数只能处处等于 $c$。

若函数\textbf{不是常数函数}，而且存在最小正周期 $T_0$，则
\[
\boxed{b-a=\frac{(2k+1)T_0}{4}\quad(k\in\mathbb Z).}
\]
即一条对称轴与一个对称中心在水平方向上的距离，是\textbf{奇数个四分之一最小正周期}。

\medskip
\textbf{\textcolor{CStatement}{使用范围务必看清}}

本结论中 $T$ 只是“一个周期”时，前两条是\textbf{保证存在的一组位置}。即使已知最小正周期，对一般抽象函数的“全部对称中心”也不能像“全部对称轴”那样不加条件地照搬；不过在函数\textbf{所有函数值都被某个固定范围夹住}（例如 $\sin x$、$\cos x$）时，可以作进一步判断，详见证明部分。
\end{statementbox}
